\documentclass{article}
\usepackage{graphicx} % Required for inserting images
\usepackage{datetime} % Add this line to include the datetime package
\usepackage{amsmath}
\usepackage{amssymb}
\usepackage{amsthm}
\usepackage{tikz}
%\usepackage{hyperref}
\usepackage{xcolor} % Required for defining custom colors

\definecolor{darkblue}{rgb}{0.0, 0.0, 0.5} % Define a darker blue color

\usepackage[colorlinks=true, linkcolor=darkblue, urlcolor=darkblue, citecolor=darkblue]{hyperref} % Use the darker blue color

% Define theorem style
\theoremstyle{plain} % Bold title, italic body
\newtheorem{theorem}{Theorem}
\newtheorem{lemma}[theorem]{Lemma}
\newtheorem{corollary}[theorem]{Corollary}

\theoremstyle{definition} % Bold title, normal body
\newtheorem{definition}[theorem]{Definition}
\newtheorem{example}[theorem]{Example}


\title{Gaussian Primes}
\author{Aaron Honculada}
\date{\monthname[\month] \the\year}

\begin{document}

\maketitle
\section{Definition}

\begin{definition}
    A Gaussian prime is an irreducible element of $\mathbb{Z}[i]$.
    Otherwise stated, it is only divisible by units, $\{\pm 1, \pm i\}$ and associates.
\end{definition}

\section{Applications}
We can use Gaussian Integers to characterize the prime numbers. We know a prime $p$ can
only be congruent to $1 \pmod{4}$ or $3 \pmod{4}$. Also recall Fermat's theorem 
on the sum of two squares, which states that a prime $p$ can be expressed as a sum of two squares
if and only if $p = 2$ or $p \equiv 1 \pmod{4}$.

\begin{theorem}
    A prime number $p$ is a Gaussian prime if and only if $p \equiv 3 \pmod{4}$.
\end{theorem}
\begin{proof}
    ($\Rightarrow$)
    For the sake of contradiction, assume $p \equiv 3 \pmod{4}$ is not a Gaussian prime.
    Then, $p$ can be factored into two non-unit Gaussian integers.
    \begin{align*}
        p &= (a + bi)(c + di)
    \end{align*}
    However when we take the norm of both sides, we get
    \[
        N(p) = N(a + bi)N(c + di)
    \]
    \[
        p^2 = N(a + bi)N(c + di)
    \]
    Since $p$ is a prime number in $\mathbb{Z}$, this implies that one of the norm, say $N(a + bi)$, is either $1$ or $p$.
    If $N(a + bi) = 1$, then $a + bi$ is a unit, which does not lead to a factorization.
    If $N(a + bi) = p$, then $p$ can be expressed as a sum of two squares, $a^2 + b^2 = p$.
    Now we show that $p \equiv 3 \pmod{4}$ is not a sum of two squares.
    By Fermat's theorem on the sum of two squares, a prime $p$ can be expressed as a sum of two squares if and only if $p = 2$ or $p \equiv 1 \pmod{4}$.
    Since we assumed $p \equiv 3 \pmod{4}$, it is not a sum of two squares.
    This contradicts the assumption that $p$ factors in $\mathbb{Z}[i]$.
    Therefore, $p$ is irreducible in $\mathbb{Z}[i]$, and must be a Gaussian prime.

    ($\Leftarrow$)
    For the sake of contradiction, assume $p \equiv 1 \pmod{4}$.
    By Fermat's theorem on the sum of two squares, $p$ can be expressed as a sum of two squares.
    \begin{align*}
        p &= a^2 + b^2
    \end{align*}
    This allows us to factor $p$ in $\mathbb{Z}[i]$ as
    \begin{align*}
        p &= (a + bi)(a - bi)
    \end{align*}
    Since neither $a + bi$ nor $a - bi$ is a unit, this is a nontrivial
    factorization of $p$ in $\mathbb{Z}[i]$, meaning $p$ is not a Gaussian prime.
    We have shown that a prime $p \equiv 1 \pmod{4}$ always factors in $\mathbb{Z}[i]$, meaning it is not a Gaussian prime.
    Therefore the only primes that are Gaussian primes are those that are $3 \pmod{4}$.
\end{proof}

\begin{corollary}
    If $p \equiv 1 \pmod{4}$, then $p$ is the norm of a Gaussian prime.
\end{corollary}
\begin{proof}
    By Fermat's theorem on the sum of two squares, $p$ can be expressed as a sum of two squares.
    \begin{align*}
        p &= a^2 + b^2
    \end{align*}
    This allows us to factor $p$ in $\mathbb{Z}[i]$ as
    \begin{align*}
        p &= (a + bi)(a - bi)
    \end{align*}
    This means that $p$ is not a prime since it factors non-trivially in $\mathbb{Z}[i]$.
    Since $N(a + bi) = p$, $a + bi$ is a Gaussian prime, as integers
    norm is a prime number in $\mathbb{Z}$.

\end{proof}


\begin{example}
    $2$ is not a Gaussian prime because $2 = (1 + i)(1 - i)$.
    However we can express $2$ as the norm of the Gaussian prime $1 + i$.
    \[
        N(1 + i) = N(1) + N(i) = 1^2 + 1^2 = 2
    \]
\end{example}

We have thus characterized all the primes in $\mathbb{Z}$ as either Gaussian primes or the norms of Gaussian primes.


\end{document}